% Direct-authority diplomatic French for printed p.130.
\SourcePage{130}

\LayerTarget{tohoku:section:1.5:p3}
\noindent d'existence maniables pour des objets ``injectifs'' et
``projectifs'' (voir N\textsuperscript{o}10). Pour éviter des redites, nous
allons d'abord étudier un procédé très général et très employé de formation de
nouvelles catégories à l'aide de diagrammes.

\LayerTarget{tohoku:section:1.6}
\noindent\textbf{1.6. Catégories de diagrammes et propriétés de permanence.}
Un \textit{schéma de diagramme} est un triple $(I,\Phi,d)$ formé de deux
ensembles $I$ et $\Phi$ et d'une application $d$ de $\Phi$ dans $I\times I$.
Les éléments de $I$ sont les \textit{sommets}, les éléments de $\Phi$ les
\textit{flèches} du diagramme, et si $\varphi$ est une flèche du diagramme,
$d(\varphi)$ est appelé sa \textit{direction}, caractérisée par
\textit{l'origine} et l'\textit{extrémité} de la flèche (ce sont donc des
sommets du schéma).
\LayerTarget{tohoku:def:composite-arrow}
Une \textit{flèche composée} d'origine $i$ et extrémité $j$ est par définition
une suite finie non vide de flèches du diagramme, l'origine de la première
étant $i$, l'extrémité de chacune étant l'origine de la suivante et
l'extrémité de la dernière étant $j$.
\LayerTarget{tohoku:def:diagram-category}
Si $\mathbf{C}$ est une catégorie, on appelle \textit{diagramme dans $\mathbf{C}$, de schéma
$\mathrm{S}$} une fonction $D$ qui associe à tout $i\in I$ un objet $D(i)\in \mathbf{C}$, et à
toute flèche $\varphi$ d'origine $i$ et extrémité $j$, un morphisme
$D(\varphi)$ de $D(i)$ dans $D(j)$. La classe de tels diagrammes sera noté
$\mathbf{C}^S$, elle sera considérée comme une catégorie, en prenant comme morphismes
de $D$ dans $D'$ une famille de morphismes
$v_i:D(i)\to D'(i)$ telle que pour toute flèche $\varphi$
d'origine $i$ et d'extrémité $j$, on ait commutativité dans le diagramme :

\LayerTarget{tohoku:diagram:diagram-morphism-square}
\[
\begin{array}{ccc}
D(i) & \stackrel{v_i}{\longrightarrow} & D'(i) \\
{D(\varphi)}\,\big\downarrow &&
  \big\downarrow\,{D'(\varphi)} \\
D(j) & \stackrel{v_j}{\longrightarrow} & D'(j)\qquad .
\end{array}
\]
Les morphismes de diagrammes se composent de façon évidente, et les axiomes
d'une catégorie se vérifient trivialement.
\LayerTarget{tohoku:def:commutative-diagram}
Si $D$ est un diagramme de schéma $\mathrm{S}$, alors pour toute flèche composée
$\varphi=(\varphi_1,\ldots,\varphi_k)$ dans $\mathrm{S}$, on définit
$D(\varphi)=D(\varphi_k)\ldots D(\varphi_1)$, c'est un morphisme de $D(i)$
dans $D(j)$ si $i$ et $j$ sont resp. l'origine et l'extrémité de $\varphi$.
$D$ est dit un \textit{diagramme commutatif} si on a
$D(\varphi)=D(\varphi')$ chaque fois que $\varphi$ et $\varphi'$ sont deux
flèches composées ayant même origine et même extrémité.
\LayerTarget{tohoku:def:commutation-relations}
Plus généralement, si $R$ est un ensemble formé de couples
$(\varphi,\varphi')$ de flèches composées ayant mêmes origines et mêmes
extrémités, et de flèches composées dont l'origine égale l'extrémité, on
considère la sous-catégorie $\mathbf{C}^{S,R}$ de $\mathbf{C}^S$ formée des diagrammes
satisfaisant aux \textit{relations de commutation}
$D(\varphi)=D(\varphi')$ pour $(\varphi,\varphi')\in R$,
$D(\varphi)={}$morphisme identique de $D(i)$ si $\varphi\in R$ a $i$ pour
origine et extrémité.

\LayerTarget{tohoku:def:additive-diagram-relations:p1}
On a à envisager encore d'autres types de commutation pour des diagrammes,
dont la nature varie avec les catégories envisagées. Ce qui suit semble
couvrir les cas les plus importants. Pour tout $(i,j)\in I\times I$, donnons
nous un ensemble $R_{ij}$ de combinaisons linéaires formelles à coéfficients
entiers de flèches composées d'origine $i$ et extrémité $j$, et, si $i=j$,
d'un élément $e_i$ auxiliaire. Si alors $D$ est un diagramme, à valeurs dans
une catégorie additive $\mathbf{C}$, on peut pour tout $L\in R_{ij}$ définir le
morphisme $D(L):D(i)\to D(j)$, en remplaçant dans l'expression de
$L$, une flèche composée $\varphi$ par $D(\varphi)$, et $e_i$ par le morphisme
identique de $D(i)$. Si on désigne par $R$ la réunion des

\TohokuContinuation{131}
